\section{固定耗时下的三类深度重构}
\label{sec:deep}
本节回应核心设计意图：\textbf{允许新算法比旧算法使用更多总运算和显存，
但让 GPU 将在线耗时压回原预算以内。}
以下不是对现有系统已经取得效果的描述，而是具备数学目标、预算变量和验证条件的研究方案。

\subsection{方案一：多模式、多场景的风险敏感采样控制}
\subsubsection{把不确定性纳入预测，而非只增加名义轨迹}
原 MPPI 对每个控制样本主要评估一条确定初始状态与既定运动模型下的轨迹。
定义不确定量
\[
\xi=(q_0,\tau_v,\kappa_{\rm slip},a_b,\alpha_b,\mathcal O_{0:H}),
\]
分别表示初始位姿、执行时间常数、侧滑参数、制动能力和未来障碍。
给定控制序列 $U_n$ 与场景 $\xi_m$，
\[
x_{n,m,\ell+1}=f(x_{n,m,\ell},u_{n,\ell},\xi_m),\qquad
J_{nm}=J(U_n;\xi_m).
\]
场景应来自经过校准的定位后验、动力学参数区间/分布及障碍预测，
不能把同一地图/标定误差对不同激光束独立抽样。
全局偏差应在一个场景内共享，才保留相关误差的物理含义。

计算张量从 $(K,H)$ 扩成 $(K,M,H)$，$n,m$ 相互独立。
同时对所有候选使用同一组 $\xi_m$，可形成 common random numbers：
比较两个候选时
\[
\Var(\widehat J_A-\widehat J_B)
=\Var(\widehat J_A)+\Var(\widehat J_B)
-2\Cov(\widehat J_A,\widehat J_B).
\]
如果共享场景使协方差为正，候选差值噪声下降。
这不保证所有问题都正相关，但避免了仅因抽到不同难度场景而错误排序。

\subsubsection{从平均风险到 CVaR 的逐步构造}
仅最小化均值
$\bar J_n=M^{-1}\sum_mJ_{nm}$
可能接受“多数很好、少数碰撞”的动作。
对风险置信水平 $\beta\in(0,1)$ 定义上尾 CVaR：
\begin{equation}
\CVaR_\beta(J)
=\min_{\eta\in\R}
\left[\eta+\frac1{1-\beta}\E(J-\eta)_+\right].
\label{eq:cvar}
\end{equation}
$\eta$ 与 $J$ 同单位，$(a)_+=\max(a,0)$。
对连续分布，在可微点
\[
\frac{d}{d\eta}\left[\eta+\frac1{1-\beta}\E(J-\eta)_+\right]
=1-\frac{\Prob(J>\eta)}{1-\beta}.
\]
令导数为零得 $\Prob(J>\eta)=1-\beta$，
故最优 $\eta$ 为 $\beta$ 分位点；代回得到上尾平均。
对有限样本则求
\[
\widehat{\CVaR}_\beta(J_n)
=\min_\eta\left[\eta+\frac1{(1-\beta)M}\sum_m(J_{nm}-\eta)_+\right],
\]
可通过排序或分位数选择精确求该经验目标，
边界样本按分数权重处理，不必硬取整数尾部数量。

新轨迹分数可定义
\begin{equation}
S_n^{\rm risk}=\bar J_n+\lambda_R\widehat{\CVaR}_\beta(J_n)+C_n^{\rm control},
\label{eq:riskcost}
\end{equation}
$\lambda_R\ge0$ 为同标度风险权重，$C_n^{\rm control}$ 是与采样方案匹配的控制修正。
然后仍可按第 4 节的熵正则推导对 $S_n^{\rm risk}$ 加权。
这改变了优化问题，目标从名义平均偏好转向尾部风险偏好。
CVaR 的标准来源见 Rockafellar--Uryasev\cite{cvarref}。
$M=16,\beta=0.9$ 时尾部只有约 $1.6$ 个样本的有效质量，
风险估计很粗；不能把很高的 $\beta$ 配很少场景当成可靠风险度量。

\subsubsection{有限场景零碰撞并非任意小概率保证}
对预先固定的某个动作，真实碰撞概率为 $p$，
$M$ 个独立同分布场景都不碰撞的概率为 $(1-p)^M$。
若用置信失误率 $\alpha$ 排除过大的 $p$，
令 $(1-p_{\rm upper})^M=\alpha$ 得
\[
p_{\rm upper}=1-\alpha^{1/M}.
\]
但控制器从 $K$ 个动作中选择，选择行为会利用场景样本。
若全部候选在抽场景前固定，按 union bound 为每个候选分配 $\alpha/K$，
可得同时有效的保守上界
\begin{equation}
p_{\rm upper}=1-(\alpha/K)^{1/M}.
\label{eq:scenario}
\end{equation}
不同候选的碰撞事件无需相互独立；需要每候选的场景抽样条件成立。
若候选依据同一场景数据继续自适应优化，则这个有限固定集合证明不再直接适用；
需独立验证场景或相应的统计学习保证。

因此 $K=1600,\alpha=0.05$ 时，
$M=32$ 的上界约 $0.277$，$M=256$ 也仍约 $0.0397$。
这些是保守统计界，说明“小批量零碰撞”不足以证明极低风险。
场景模型失配时即使统计界形式正确也不能保证实物风险；
对制动等关键量可额外采用物理可信的区间最坏情形。
GPU 能增加检查覆盖，不能消除模型校准这一条件。

\subsubsection{多模式输出：避免把左绕和右绕平均到障碍上}
可行控制集合通常非凸。
若左绕序列 $U_L$ 与右绕序列 $U_R$ 都安全，
其平均 $(U_L+U_R)/2$ 可能直行穿过障碍。
因此“每条样本都安全”不保证 MPPI 加权均值安全。

建议按停车、保持路径、左绕、右绕等行为族分别进行采样和权重更新，
得到各族候选 $\bar U_a$；
重新评估 $\bar U_a$ 本身及实际第一步过渡，
再从通过验证者中选择。
模式切换加入滞回，避免噪声导致左右来回切换。
重新评估是必要计算，不可因 GPU 已检查样本就省略。
这项重构改变的不仅是样本数，还改变了输出决策结构。

\subsubsection{在原预算内如何选轨迹数、场景数与步数}
基线 $KH=89600$。
一个\emph{预算候选}为 $K=800,M=8,H=56$、两轮更新，
状态展开数为 $716800$，约基线 8 倍。
如果式 \eqref{eq:expansion} 的实测允许扩张量至少为 8，
且新增风险归约、地图访问、最终候选复验也包含在时间中，
它才有机会满足同耗时要求。
这不是已经验证能实时的配置。

单轮同时保存 $K M H$ 个三维 float 位姿需
$800\times8\times56\times3\times4=4.3008$ MB。
若每场景维护更多状态和代价，内存相应增加。
可在轨迹展开时累积可加代价，减少全轨迹存储；
路径对齐等需要历史信息时则不能无条件流式丢弃。
若地图查询是主要瓶颈，新增场景未必按算力比例加速。
可采用“低成本全体筛选 + 高风险/高质量候选增加场景”的分层预算，
但筛选误差和自适应选择的统计偏差必须独立评估。

\subsection{方案二：显式位姿概率网格与层次多假设定位}
\subsubsection{用显存换分布覆盖，而非只保存有限粒子}
在离散状态空间
$\mathcal X=\{(i,j,k):0\le i<N_x,0\le j<N_y,0\le k<N_\theta\}$
上保存概率质量 $b_t(i,j,k)$，
$\sum_{i,j,k}b_t=1$。
给定离散运动转移矩阵 $P(x|x')$，
Bayes 递推为
\begin{align}
b_t^-(x)&=\sum_{x'}P(x|x')b_{t-1}^+(x'),\\
b_t^+(x)&=\frac{L_t(x)b_t^-(x)}{\sum_yL_t(y)b_t^-(y)}.
\label{eq:gridbayes}
\end{align}
这是选定有限状态模型的完整分布递推，无粒子重采样丢失模式问题；
但有网格量化、截断域和模型误差。
精确离散递推不等于精确连续定位。

若直接使用稠密转移矩阵，复杂度为状态数平方，显然不适合。
重构关键在利用运动核的局部性和按方向的平移不变结构。

\subsubsection{由全向运动模型得到卷积}
对旧方向层 $k'$，把平移噪声及平均里程计位移形成二维离散概率核
$G_{k'}(\Delta i,\Delta j)$。
它可包含随方向旋转的各向异性协方差与非零均值。
令角度增量核为 $A(k|k')$。
在“给定旧方向后平移与角增量独立、空间中运动模型平移不变”的假设下，
\[
P(i,j,k|i',j',k')
=A(k|k')G_{k'}(i-i',j-j').
\]
代入式 \eqref{eq:gridbayes}，交换有限求和：
\begin{align}
c_{k'}(i,j)&=
\sum_{i',j'}G_{k'}(i-i',j-j')b_{t-1}^+(i',j',k'),\\
b_t^-(i,j,k)&=\sum_{k'}A(k|k')c_{k'}(i,j).
\label{eq:convpredict}
\end{align}
第一步是每个方向的二维卷积，第二步是方向混合。
若角度核有效宽度为 $B_\theta$，直接方向混合工作
$O(GB_\theta N_\theta)$；平移卷积可局部 stencil 或 FFT。
FFT 用
$\operatorname{FFT}(G*b)=\operatorname{FFT}(G)\operatorname{FFT}(b)$，
每层工作约 $O(G\log G)$。
线性卷积必须零填充或保持足够 halo，
否则 FFT 的循环边界会把一侧概率错误送到另一侧。

若截断高斯运动核，必须记录丢弃尾质量并重新定义近似模型。
边界离开局部窗口的概率不能凭空归一化后假装没有发生；
可转入外层粗地图或保留“窗口外”质量。
测量更新可能放大先验近似误差，所以仅知道转移截断很小，
也不能无条件保证后验误差同样小。

\subsubsection{把全网格激光评分改写成相关运算}
当前 beam-skip 关闭且每束使用同一评分函数。
定义地图函数
\[
g(x)=\log\left[a\exp(-D_{\mathcal M}(x)^2/(2\sigma^2))+b_0\right].
\]
在固定方向 $\theta_k$、固定平移网格及一致的端点量化规则下，
令 $\Delta_{kb}$ 为第 $b$ 束对应的整数端点偏移。
创建稀疏扫描模板
$S_k(\Delta)=\sum_b\ind\{\Delta=\Delta_{kb}\}$；
重复落格的束以计数保留。
则该方向整幅位姿图的对数观测评分
\begin{align}
\ell_k(p)
&=\sum_bg(p+\Delta_{kb})\\
&=\sum_\Delta S_k(\Delta)g(p+\Delta).
\label{eq:senscorr}
\end{align}
第二行正是互相关，可用 FFT 实现。
在支持条件成立时，原 $O(N_\theta GB)$ 的点束评分，
转成约 $O(N_\theta G_{\rm pad}\log G_{\rm pad})$；
$G_{\rm pad}$ 是为线性相关扩大的图。
图的 FFT 可以复用，扫描模板随扫描变化。
这是一项比“增加线程”更深入的算法结构变化。

限制也明确：端点量化必须与原 grid 坐标一致；
逐粒子束跳过、位置相关束权、变化的可见性模型可能破坏简单相关形式；
越界必须用原截断距离对应的常数评分填充，
不能默认零或周期地图。
数值 FFT 会有误差，不能仅凭代数恒等式宣称逐位相同。
模板旋转的量化若改变端点，需单独分析。
从成本角度，$B=120$ 与 $\log G$ 的比较只是渐近判断，
常数、填充和 FFT 批量大小仍需测量。

\subsubsection{规模算例及合理的层次结构}
\begin{center}\small
\begin{tabular}{lrrr}
\toprule 位姿网格 & 状态数 & 单 float 场内存 & 直接 120 束评分量\\\midrule
局部 $4$ m，$4$ cm，64 方向 & $640000$ & $2.56$ MB & $76800000$\\
局部 $4$ m，$1$ cm，64 方向 & $10240000$ & $40.96$ MB & $1228800000$\\
全局 $20$ m，$1$ cm，64 方向 & $256000000$ & $1024$ MB & $30720000000$\\
\bottomrule
\end{tabular}
\end{center}
内存仅为一个标量场，未计预测/后验双缓冲、对数似然和 FFT 复数临时区。
最后一行说明显存能容纳不代表扫描周期能完成计算。
因此不建议第一步就做全地图毫米级三维稠密 Bayes。

更合理路线：全局较粗分布维持所有有意义模式，
对若干高质量且空间分离的模式建立局部细网格，
局部用连续匹配或粒子更新。
粗细层应是互斥分区的概率质量，或显式混合：
局部层权重由粗层分配，不能把同一状态质量在两层重复加和。
细化区域扩张、合并和删除时保持总质量；
不应只保留 top-1 后将其权重硬设为 1。
GPU 的额外存储在这里用于\emph{保留歧义}，
不是制造一个数值上很尖却可能错误的后验。

\subsection{方案三：预计算配置空间支撑更大动作图及主动消歧}
\subsubsection{空间换时间，让在线搜索能覆盖更多行为}
式 \eqref{eq:cspace} 把地图与足迹组合编译成位置--方向可行表。
当 $Q$ 次在线查询共享同一地图，原暴力成本约
$Q\,c\,N_oN_v$，
新成本约
$T_{\rm build}+Q\,t_{\rm lookup}+Q_{\rm near}t_{\rm exact}$。
需同时满足
\begin{equation}
T_{\rm build}+Q\,t_{\rm lookup}+Q_{\rm near}t_{\rm exact}
<Q\,cN_oN_v
\end{equation}
以及单周期更新峰值预算。
节省下的查询时间可用于增加航向/横移候选、多个全局通路、
更多控制情景或更细时间步长。
“是否值得预计算”取决于共享查询数 $Q$，
而不取决于单张表能否放进显存。

静态图和动态障碍可分层：静态配置空间预计算，
动态局部障碍每周期更新较小区域，
在线碰撞是两层碰撞结果的 OR。
若风险情景包含不同障碍轨迹，可以共享静态层，
按情景建立动态层或对动态障碍精查。
不得因为静态地图通过就忽略扫描中新出现的障碍。

\subsubsection{主动感知：让控制帮助定位解决歧义}
如果定位后验同时有多个模式，仅做更多计算可能仍无法区分重复走廊。
机器人可以选择能看到岔口/非对称结构的微小动作，
用下一次真实观测获得新信息。
这时额外 GPU 预算用于“候选动作 $\times$ 位姿假设 $\times$ 预测观测”。

设离散模式随机变量为 $A$，权重 $\pi_a$；
候选动作 $u$ 后的未来激光特征为 $Z$，模型密度 $p(z|a,u)$。
混合观测密度
$p(z|u)=\sum_a\pi_ap(z|a,u)$。
由 Bayes，
$\pi_a(z,u)=\pi_ap(z|a,u)/p(z|u)$。
先验熵 $H(A)=-\sum_a\pi_a\log\pi_a$，
期望后验熵为 $\int p(z|u)[-\sum_a\pi_a(z,u)\log\pi_a(z,u)]\,dz$。
相减并代入 Bayes，化简得
\begin{equation}
I(A;Z|u)=
\sum_a\pi_a\int p(z|a,u)
\log\frac{p(z|a,u)}{p(z|u)}\,dz.
\label{eq:infogain}
\end{equation}
可以把 $-\lambda_I I(A;Z|u)$ 加到任务代价，同时约束运动风险与任务进度。

直观地，两模式若预测同一测量高斯
$\mathcal N(\mu_a,\sigma_z^2)$，
其 KL 散度为
\[
D_{\rm KL}(\mathcal N(\mu_1,\sigma_z^2)\|
\mathcal N(\mu_2,\sigma_z^2))
=\frac{(\mu_1-\mu_2)^2}{2\sigma_z^2}.
\]
证明是将两正态对数密度相减，
对第一分布取期望：
$\E[(Z-\mu_2)^2-(Z-\mu_1)^2]=(\mu_1-\mu_2)^2$。
这只是两模式可区分度，不等于式 \eqref{eq:infogain} 的混合互信息；
但说明动作若扩大预测观测差异，有潜力增加消歧信息。
若所有可执行动作下预测都相同，任何计算能力也无法产生信息。

该方向最能体现“用算力提升算法能力”，同时开发与建模成本最高。
应放在可靠多假设定位和场景控制之后，
先验证预测观测模型，再引入主动动作。
同一激光既用于定位又用于场景生成时，应使用一个联合后验，
不能把重复使用的证据当作独立观测再次相乘。
